تفاوت‌های کلیدی بین \lr{image enhancement} و \lr{image restoration} به صورت خلاصه:

1) هدف: بهبود دیداری و ادراکی در \lr{enhancement} در برابر بازسازی نزدیک به تصویر اصلی در \lr{restoration}.

2) ماهیت مسئله: \lr{enhancement} معمولاً ذهنی و وابسته به کاربرد است؛ \lr{restoration} عینی و مدل‌محور است.

3) اطلاعات لازم: \lr{enhancement} غالباً بدون مدل تخریب؛ \lr{restoration} نیازمند مدل نویز/تخریب است.

4) معیار ارزیابی: \lr{enhancement} با معیارهای ادراکی/کیفی؛ \lr{restoration} با معیارهای کمی مانند \lr{MSE} یا \lr{PSNR}.

5) خروجی: \lr{enhancement} ممکن است جزئیات را اغراق کند؛ \lr{restoration} تلاش می‌کند اعوجاج‌ها را حذف و وفاداری را حفظ کند.